\title{Direct Preference Optimization: Your Language Model is Secretly a Reward Model}

\begin{abstract}
Although large-scale unsupervised language models (LMs) acquire extensive world knowledge and some reasoning capabilities, precisely directing their behavior remains challenging due to the fully unsupervised nature of their training. Current techniques for achieving such control gather human annotations on the relative quality of model outputs and then fine-tune the unsupervised LM to align with these preferences, often employing reinforcement learning from human feedback (RLHF). Nevertheless, RLHF is a complicated and frequently unstable process, involving first training a reward model that captures human preferences, and subsequently fine-tuning the large unsupervised LM via reinforcement learning to maximize this inferred reward while avoiding excessive deviation from the original model. In this work, we propose a novel parameterization of the reward model within RLHF that permits deriving the corresponding optimal policy in closed form, enabling us to address the standard RLHF problem using only a straightforward classification loss. The resulting method, which we name \textit{Direct Preference Optimization} (DPO), is stable, effective, and computationally efficient, removing the necessity for sampling from the LM during fine-tuning or extensive hyperparameter optimization. Our experimental results demonstrate that DPO can fine-tune LMs to align with human preferences as well as or better than existing approaches. Importantly, fine-tuning with DPO outperforms PPO-based RLHF in controlling the sentiment of outputs, and matches or surpasses response quality in summarization and single-turn dialogue tasks, all while being considerably simpler to implement and train.
\end{abstract}

\section{Introduction}  
Large-scale unsupervised language models (LMs) trained on extensive datasets develop surprisingly advanced capabilities~\citep{chowdhery2022palm, brown2020language, touvron2023llama,bubeck2023sparks}. However, these models learn from human-generated data that reflects a diverse range of objectives, priorities, and expertise. Some of these objectives and skillsets may not be desirable to replicate; for instance, while it is beneficial for an AI coding assistant to \textit{recognize} common programming errors to fix them, when generating code, we want to steer the model towards the (possibly rare) instances of high-quality coding found within its training data. Likewise, we may want a language model to be \textit{aware} of a misconception held by about 50\% of people, but we certainly do not want the model to assert this misconception as true in half of all responses concerning it! Put differently, carefully choosing the model's \emph{preferred responses and behaviors} from its broad \textit{knowledge and capabilities} is essential for developing AI systems that are reliable, effective, and controllable \citep{ouyang2022training}. Although current approaches generally guide LMs to align with human preferences through reinforcement learning (RL), we will demonstrate that the RL objective employed by these methods can be exactly optimized using a straightforward binary cross-entropy loss, thereby greatly simplifying the preference learning process.

Broadly, existing techniques impart desired behaviors into a language model by leveraging carefully curated human preference datasets that represent behaviors regarded as safe and beneficial. This preference learning phase follows an initial stage of large-scale unsupervised pre-training on massive text corpora. While the simplest method for preference learning is supervised fine-tuning on human demonstrations of high-quality responses, the most effective approaches utilize reinforcement learning from human (or AI) feedback (RLHF/RLAIF; \citep{christiano2017deep,bai2022constitutional}). RLHF approaches involve fitting a reward model to a dataset of human preferences and then applying RL to optimize a language model policy to generate responses with high predicted reward, while ensuring the model does not diverge excessively from the original. Although RLHF yields models with strong conversational and coding skills, the RLHF framework is significantly more complex than supervised learning, requiring training multiple language models and sampling from the model policy during training, which leads to substantial computational demands.

In this work, we demonstrate how to optimize a language model directly to align with human preferences, without resorting to explicit reward modeling or reinforcement learning. We introduce \textit{Direct Preference Optimization} (DPO), an algorithm that implicitly optimizes the same objective as current RLHF methods (reward maximization subject to a KL-divergence constraint) while being easy to implement and simple to train. Conceptually, the DPO update raises the relative log probability of preferred responses compared to dispreferred ones, incorporating a dynamic, per-example importance weighting that mitigates model collapse observed with a straightforward probability ratio objective. Similar to existing approaches, DPO depends on a theoretical preference model (e.g., the Bradley-Terry model; \cite{bradley1952rankanalysis}) which evaluates how well a reward function corresponds with observed preference data. However, whereas current methods use the preference model to create a preference loss for training a reward model, and subsequently train a policy to optimize this learned reward model, DPO performs a change of variables to define the preference loss directly in terms of the policy. Thus, given a dataset containing human preferences over model outputs, DPO can optimize a policy through a simple binary cross entropy loss, yielding the optimal policy for an implicit reward function fitted to the preference data.

Our primary contribution is Direct Preference Optimization (DPO), a straightforward, reinforcement learning–free algorithm for training language models based on preferences. Our experiments indicate that DPO is at least as effective as existing approaches, including PPO-based RLHF, in learning from preferences on tasks such as sentiment modulation, summarization, and dialogue, using language models with up to 6 billion parameters.

\section{Related Work}

Self-supervised language models of progressively larger scale can perform certain tasks either zero-shot \citep{radford2019language} or with few-shot prompting \citep{gpt3,megatron,chowdhery2022palm}. Nevertheless, their effectiveness on downstream tasks and alignment with user intent can be substantially enhanced by fine-tuning on datasets containing instructions paired with human-generated completions \citep{mishra-etal-2022-cross,sanh2022multitask,chung2022scaling,thoppilan2022lamda}. This process, known as `instruction-tuning,' allows LLMs to generalize beyond the specific instructions in the tuning set and generally improve their usability \citep{chung2022scaling}. Despite the success of instruction tuning, \textit{relative} human judgments regarding response quality are frequently easier to obtain than expert demonstrations. Consequently, later studies have fine-tuned LLMs using datasets composed of human preferences, which has enhanced capabilities in areas like translation \citep{kreutzer-etal-2018-reliability}, summarization \citep{stiennon2022learning,ziegler2020finetuning}, storytelling \citep{ziegler2020finetuning}, and instruction-following \citep{ouyang2022training,ramamurthy2023is}. These approaches first train a neural network reward function to align with the preference dataset under a preference model such as the Bradley-Terry model \citep{bradley1952rankanalysis}, then fine-tune a language model to maximize this reward via reinforcement learning algorithms, typically REINFORCE \citep{williams1992reinforce}, proximal policy optimization (PPO; \cite{schulman2017proximal}), or their variants \citep{ramamurthy2023is}. A closely related research thread uses LLMs that have been fine-tuned with human feedback for instruction-following to create additional synthetic preference datasets focusing on specific attributes like safety or harmlessness \citep{bai2022constitutional}, relying solely on weak human supervision through textual rubrics for the LLM's annotations. These techniques signify a merging of two research domains: one concerning training language models with reinforcement learning for various objectives \citep{Ranzato2015SequenceLT,paulus2018a,wu2018learning}, and another focusing on general methods to learn from human preferences \citep{christiano2017deep,kupcsik2018learning}. Although leveraging relative human preferences is attractive, applying reinforcement learning to fine-tune large language models remains a significant practical challenge; this work proposes a theoretically grounded method to optimize relative preferences without relying on RL.

Outside the realm of language, the problem of learning policies from preferences has been investigated within both bandit and reinforcement learning frameworks, with numerous methodologies proposed. Contextual bandit learning that relies on preferences or rankings of actions instead of rewards is referred to as a contextual dueling bandit (CDB; \cite{yue2012karmed,dudik2015contextual}). When absolute rewards are unavailable, the theoretical framework for CDBs replaces the concept of an optimal policy with the \textit{von Neumann winner}, which is a policy that has an expected win rate of at least 50\% against \textit{any} alternative policy \citep{dudik2015contextual}. Nevertheless, in the CDB context, preference labels are provided online, whereas in learning from human preferences, the typical scenario involves training from a static batch of offline preference-annotated pairs of actions \citep{yan2022human}. Likewise, \textit{preference-based RL} (PbRL) learns from binary preferences that are generated by an \textit{unknown} `scoring' function rather than explicit reward signals \citep{BusaFekete2014,ruiz2023dueling}. There exist various algorithms for PbRL, including those capable of reusing off-policy preference data, but these generally entail first explicitly estimating the latent scoring function (i.e., the reward model), followed by its optimization \citep{jain2013learning,BusaFekete2014,christiano2017deep,sadigh2017active,kupcsik2018learning}. In contrast, we propose a single-stage policy learning method that directly optimizes the policy to align with the observed preferences.

\section{Preliminaries}\label{section:prelims}

We revisit the RLHF pipeline as described in \citeauthor{ziegler2020finetuning} (and subsequently by \citep{stiennon2022learning, bai2022training, ouyang2022training}). This pipeline typically involves three steps: 1) supervised fine-tuning (SFT); 2) preference sampling and reward learning; and 3) reinforcement learning optimization.

\textbf{SFT}: RLHF generally starts by fine-tuning a pre-trained language model (LM) using supervised learning on high-quality datasets relevant to the downstream tasks of interest (such as dialogue or summarization), resulting in a model $\pi^\text{SFT}$.

\textbf{Reward Modelling Phase}: In the second step, the SFT model is prompted with inputs $x$ to generate pairs of responses $(y_1, y_2) \sim \pi^\text{SFT}(y \mid x)$. These response pairs are then shown to human annotators who indicate their preferences, denoted as $y_w \succ y_l \mid x$, where $y_w$ and $y_l$ represent the preferred and less preferred completions from the pair $(y_1, y_2)$, respectively. These preferences are assumed to be derived from some latent reward function $r^*(y, x)$, which is not directly observable. Various methods exist to model these preferences; a common choice is the Bradley-Terry (BT) model \cite{bradley1952rankanalysis}, although more general Plackett-Luce ranking models \citep{plackett1975analysis, luce2012individual} can also be integrated into the framework when multiple ranked responses are available. According to the BT model, the human preference distribution $p^*$ can be expressed as:

\begin{equation}\label{eq:bradley-terry}
    p^*(y_1\succ y_2 \mid x)=\frac{\exp\left(r^*(x, y_1)\right)}{\exp\left(r^*(x, y_1)\right) + \exp\left(r^*(x, y_2)\right)}.
\end{equation}

Given a fixed dataset of comparisons $\mathcal{D}=\bigl\{x^{(i)}, y_w^{(i)}, y_l^{(i)}\bigr\}_{i=1}^N$ drawn from $p^*$, one can parameterize a reward model $r_{\phi}(x, y)$ and estimate its parameters through maximum likelihood estimation. By casting the problem as a binary classification task, the negative log-likelihood loss is expressed as:

\begin{equation}\label{eq:reward_model}
    \mathcal{L}_R(r_{\phi}, \mathcal{D}) = -\mathbb{E}_{(x, y_w, y_l)\sim \mathcal{D}}\bigl[\log \sigma(r_{\phi}(x, y_w)- r_{\phi}(x, y_l))\bigr]
\end{equation}

where $\sigma$ denotes the logistic function. In the setting of language models, the network $r_{\phi}(x, y)$ is often initialized from the supervised fine-tuning (SFT) model $\pi^\text{SFT}(y \mid x)$, augmented by a linear layer atop the final transformer layer to output a scalar reward prediction \cite{ziegler2020finetuning}. To reduce variance in the reward signal, previous works normalize the rewards such that $\mathbb{E}_{x,y\sim \mathcal{D}}\left[r_\phi(x, y)\right] = 0$ for all $x$.

\textbf{RL Fine-Tuning Phase}: In the reinforcement learning stage, the learned reward function provides feedback to the language model. Following earlier studies~\citep{jaques2017sequence, jaques2020human}, the optimization objective is formulated as

\begin{equation}\label{eq:RL}
\max_{\pi_{\theta}}  \mathbb{E}_{x\sim \mathcal{D}, y\sim \pi_{\theta}(y \mid x)}\bigl[r_{\phi}(x, y)\bigr] - \beta\mathbb{D}_{\textrm{KL}}\bigl[\pi_{\theta}(y\mid x)\mid \mid \pi_\text{ref}(y\mid x)\bigr],
\end{equation}

where $\beta$ is a hyperparameter that regulates the divergence from the base reference policy $\pi_\text{ref}$, which is the initial SFT model $\pi^\text{SFT}$. In practice, the language model policy $\pi_\theta$ is also initialized with $\pi^\text{SFT}$. This regularization is critical to prevent the model from straying too far from the distribution where the reward model is reliable, while also preserving diversity in generated outputs and avoiding collapse to a limited set of high-reward responses. Since language generation is discrete, this objective is not differentiable and is generally optimized via reinforcement learning. The common method \citep{ziegler2020finetuning, stiennon2022learning, bai2022training, ouyang2022training} is to define the reward as ${r(x, y) = r_{\phi}(x, y) -\beta (\log \pi_{\theta}(y\mid x) - \log \pi_\text{ref}(y\mid x))}$ and maximize it using PPO \cite{schulman2017proximal}.

\section{Direct Preference Optimization}\label{sec:DPO}

Due to the difficulties in applying reinforcement learning algorithms to large-scale tasks such as fine-tuning language models, our aim is to develop a straightforward method for policy optimization that directly uses preference data. In contrast to traditional RLHF approaches, which first learn a reward function and subsequently optimize it with RL, our method exploits a specific parameterization of the reward model that allows direct derivation of the optimal policy in closed form, circumventing the need for an RL training loop. As detailed below, the central idea is to utilize an analytical relationship between reward functions and their corresponding optimal policies, enabling us to convert a loss defined over reward functions into one over policies. This change-of-variables technique obviates the necessity to explicitly fit a separate reward model, while still operating under established models of human preferences like the Bradley-Terry model. Essentially, the policy network simultaneously encodes the language model and the

\textbf{Deriving the DPO objective.} We begin with the same reinforcement learning objective as previous studies, Eq.~\ref{eq:RL}, defined with a general reward function $r$. Following earlier work~\citep{peters2007reinforcement, peng2019advantage, korbak2022reinforcement, go2023aligning}, it is straightforward to demonstrate that the optimal solution to the KL-constrained reward maximization problem in Eq.~\ref{eq:RL} has the form:

\begin{equation}\label{eq:op_policy}
    \pi_r(y\mid x) = \frac{1}{Z(x)}\pi_\text{ref}(y\mid x)\exp\left(\frac{1}{\beta}r(x, y)\right),
\end{equation}

where the partition function $Z(x) = \sum_{y} \pi_\text{ref}(y\mid x)\exp\left(\frac{1}{\beta}r(x, y)\right)$. A full derivation can be found in Appendix \ref{app:derivation1}. Even when employing the MLE estimate $r_{\phi}$ for the true reward function $r^*$, calculating the partition function $Z(x)$ remains computationally intensive~\citep{korbak2022reinforcement, go2023aligning}, which complicates practical use of this formulation. Nonetheless, we can rearrange Eq.~\ref{eq:op_policy} to express the reward function in terms of its corresponding optimal policy $\pi_r$, the reference policy $\pi_\text{ref}$, and the unknown partition function $Z(\cdot)$. Specifically, by taking the logarithm on both sides of Eq.~\ref{eq:op_policy} and performing some algebraic manipulation, we obtain:

\begin{equation}\label{eq:main_eq}
    r(x,y) = \beta \log \frac{\pi_r(y\mid x)}{\pi_\text{ref}(y\mid x)} + \beta \log Z(x).
\end{equation}

This reparameterization can be applied to the true reward $r^*$ and its corresponding optimal policy $\pi^*$. Conveniently, the Bradley-Terry model depends solely on the difference in rewards between two completions, i.e., ${p^*(y_1 \succ y_2 \mid x) = \sigma(r^*(x, y_1) - r^*(x, y_2))}$. Substituting the reparameterization from Eq.~\ref{eq:main_eq} for $r^*(x,y)$ into the preference model Eq.~\ref{eq:bradley-terry}, the partition function cancels out, allowing us to represent the human preference probability exclusively in terms of the optimal policy $\pi^*$ and reference policy $\pi_\text{ref}$. Therefore, the optimal RLHF policy $\pi^*$ under the Bradley-Terry model satisfies:

\begin{equation}\label{eq:objective}
    p^*(y_1\succ y_2 \mid x)=\frac{1}{1 + \exp\left(\beta \log \frac{\pi^*(y_2\mid x)}{\pi_\text{ref}(y_2\mid x)} - \beta \log \frac{\pi^*(y_1\mid x)}{\pi_\text{ref}(y_1\mid x)}\right)}.
\end{equation}

The full derivation is provided in Appendix~\ref{app:derivation2}. Although Eq.~\ref{eq:objective} is based on the Bradley-Terry model, analogous derivations can be done for the more general Plackett-Luce models~\citep{plackett1975analysis, luce2012individual}, as detailed in Appendix~\ref{app:plackett_luce_models}.

Having expressed the human preference data probability in terms of the optimal policy rather than the reward model, we can now define a maximum likelihood objective for a parameterized policy $\pi_\theta$. Mirroring the reward modeling approach (Eq.~\ref{eq:reward_model}), the policy objective becomes:

\begin{equation}\label{eq:optimum_model}
    \mathcal{L}_\text{DPO}(\pi_{\theta}; \pi_\text{ref}) = -\mathbb{E}_{(x, y_w, y_l)\sim \mathcal{D}}\left[\log \sigma \left(\beta \log \frac{\pi_{\theta}(y_w\mid x)}{\pi_\text{ref}(y_w\mid x)} - \beta \log \frac{\pi_{\

In this manner, we fit an implicit reward by employing an alternative parameterization, whose optimal policy corresponds directly to $\pi_\theta$. Additionally, since our approach is analogous to fitting a reparameterized Bradley-Terry model, it benefits from certain theoretical guarantees, including consistency under appropriate assumptions about the preference data distribution \cite{bong2022generalized}. We further elaborate on the theoretical aspects of DPO and its relationship to other studies in Section~\ref{sec:theory}.

\textbf{What is the effect of the DPO update?} To gain a mechanistic insight into DPO, it is helpful to examine the gradient of the loss function $\mathcal{L}_\text{DPO}$. The gradient with respect to the parameters $\theta$ can be expressed as:

\begin{multline*}\label{eq:gradient}
    \nabla_\theta \mathcal{L}_\text{DPO}(\pi_\theta;\pi_\text{ref}) = \\ -\beta\mathbb{E}_{(x, y_w, y_l) \sim \mathcal{D}} \bigg[\underbrace{\sigma(\hat{r}_\theta(x, y_l) - \hat{r}_\theta (x, y_w))}_\text{higher weight when reward estimate is wrong}\bigg[\underbrace{\nabla_\theta\log \pi(y_w \mid x)}_\text{increase likelihood of $y_w$} - \underbrace{\nabla_\theta\log\pi(y_l \mid x)}_\text{decrease likelihood of $y_l$}\bigg]\bigg],
\end{multline*}

where $\hat{r}_\theta(x, y) = \beta \log \frac{\pi_\theta(y \mid x)}{\pi_\text{ref}(y \mid x)}$ denotes the reward implicitly defined by the language model $\pi_\theta$ and the reference model $\pi_\text{ref}$ (detailed further in Section~\ref{sec:theory}). Intuitively, the gradient of $\mathcal{L}_\text{DPO}$ acts to raise the probability of preferred completions $y_w$ while lowering that of dispreferred completions $y_l$. Crucially, the training samples are weighted by the extent to which the implicit reward model $\hat{r}_\theta$ mistakenly ranks the dispreferred completions higher, scaled by $\beta$, reflecting how inaccurately the implicit reward model orders these completions and incorporating the strength of the KL divergence constraint. Our experimental results highlight the significance of this weighting, as a straightforward variant lacking the weighting factor can lead to degeneration of the language model (see Appendix Table~\ref{tab:unlikelihood_generations}).

\textbf{DPO outline.} The standard DPO process involves the following steps: 1) For each prompt $x$, generate completions $y_1, y_2 \sim \pi_\text{ref}(\cdot \mid x)$, then use human preference labels to create an offline preference dataset $\mathcal{D} = \{x^{(i)}, y_w^{(i)}, y_l^{(i)}\}_{i=1}^N$, and 2) train the language model $\pi_\theta$ by minimizing $\mathcal{L}_\text{DPO}$ given $\pi_\text{ref}$, $\mathcal{D}$, and the chosen $\beta$. In practical scenarios, it is preferable to utilize publicly accessible preference datasets instead of producing samples and collecting human preferences anew. Since these preference datasets are generated using $\pi^\text{SFT}$, we typically set $\pi_\text{ref} = \pi^\text{SFT}$ when available. If $\pi^\text{SFT}$ is not accessible, we initialize $\pi_\text{ref}$ by maximizing the likelihood over preferred completions ${(x, y_w)}$, i.e., ${\pi_\text{ref} = \argmax_{\pi} \mathbb{E}_{x, y_w \sim \mathcal{D}}\left[\log \pi(y_w \mid x)\right]}$. This approach helps reduce the distributional discrepancy between the unknown true reference distribution and the $\pi_\text{ref}$ used in DPO. Additional details on implementation and hyperparameters are provided in Appendix~\ref{app:implementation}.

\section{Theoretical Analysis of DPO} In this section, we further interpret the DPO method, present theoretical justification, and connect the benefits of DPO to challenges associated with actor-critic algorithms employed in RLHF (such as PPO~\cite{schulman2017proximal}).

\label{sec:theory}

\subsection{Your Language Model Is Secretly a Reward Model} DPO circumvents the need to explicitly fit a reward model and to perform reinforcement learning by optimizing a single maximum likelihood objective. Note that the optimization target in Eq.~\ref{eq:main_eq} corresponds to a Bradley-Terry model parameterized by a reward function $r^*(x, y) = \beta \log\frac{\pi^*_\theta(y \mid x)}{\pi_\text{ref}(y \mid x)}$. Optimizing our parametric model $\pi_\theta$ is thus equivalent to optimizing the reward model in Eq.~\ref{eq:reward_model} under this change of variables. In this section, we develop the theoretical framework underlying this reparameterization, demonstrate that it does not limit the class of learnable reward models, and show that it permits exact recovery of the optimal policy. We start by defining an equivalence relation between reward functions.

\begin{definition}
We define two reward functions $r(x, y)$ and $r'(x, y)$ to be equivalent if and only if there exists a function $f$ such that ${r(x, y) - r'(x, y) = f(x)}$.     
\end{definition}

It is straightforward to verify that this relation is an equivalence relation, which divides the set of reward functions into distinct equivalence classes. We can then present the following two lemmas:

\begin{lemma}\label{lemma:same_prefrence} Within the Plackett-Luce model, and specifically the Bradley-Terry framework, any two reward functions belonging to the same equivalence class generate identical preference distributions.
\end{lemma}

\begin{lemma}\label{lemma:same_policy}
    Any two reward functions from the same equivalence class yield the same optimal policy in the constrained reinforcement learning setting.
\end{lemma}

The proofs are direct and provided in Appendix \ref{app:lemma1}. The first lemma reflects a known identifiability issue inherent in the Plackett-Luce family of models \cite{plackett1975analysis}. Due to this ambiguity, it is usually necessary to apply extra identifiability constraints to establish guarantees for the MLE estimators derived from Eq. \ref{eq:reward_model} \cite{bong2022generalized}. The second lemma confirms that all reward functions within an equivalence class correspond to the same optimal policy, implying that for our ultimate goal, it suffices to identify any representative reward function from that class. We establish the following theorem in Appendix~\ref{app:thm1}:

\begin{theorem}\label{thm:main}
    Under reasonable assumptions, every equivalence class of reward functions consistent with the Plackett-Luce model (including the Bradley-Terry case) can be expressed in the form ${r(x, y) = \beta \log \frac{\pi(y\mid x)}{\pi_\text{ref}(y\mid x)}}$ for some policy model $\pi(y \mid x)$ and a fixed reference policy $\pi_\text{ref}(y \mid x)$.
\end{theorem}

\begin{sproof}
    Take an arbitrary reward function $r(x, y)$ that induces the corresponding optimal policy $\pi_r(y \mid x)$ as defined in Eq. \ref{eq:op_policy}. We will demonstrate that there exists a reward function within the equivalence class of $r$ that admits the given reparameterization. Define the mapping $f$ as  
\begin{equation}
    f(r; \pi_\text{ref}, \beta)(x, y) = r(x, y) - \beta \log \sum_{y} \pi_\text{ref}(y \mid x) \exp\left(\frac{1}{\beta} r(x, y)\right).
\end{equation}
This operator $f$ normalizes the reward by subtracting the logarithm of the partition function associated with $\pi_r$. Since the normalization term depends only on the prefix $x$, $f(r; \pi_\text{ref}, \beta)(x, y)$ remains within the equivalence class of the original reward $r(x, y)$. Substituting $r$ with the right-hand side of Eq.~\ref{eq:main_eq} (valid for any reward function), we obtain 
$f(r; \pi_\text{ref}, \beta)(x, y) = \beta \log \frac{\pi_r(y \mid x)}{\pi_\text{ref}(y \mid x)}$. 
Hence, the operator $f$ produces a reward function from the equivalence class of $r$ in the desired form, showing that the proposed reparameterization does not impose any loss of generality on the reward model.
\end{sproof}

An alternative interpretation of Theorem~\ref{thm:main} is that it precisely identifies which reward function within each equivalence class the DPO reparameterization chooses, specifically the reward function that satisfies:

\begin{equation}\label{eq:lag_p}
     \sum_{y}\underbrace{\pi_\text{ref}(y\mid x)\exp\left(\frac{1}{\beta}r(x, y)\right)}_{=\pi(y\mid x)\text{, using Thm.~\ref{thm:main} reparam.}} = 1,
\end{equation}

meaning that $\pi(y\mid x)$ is a valid probability distribution (all probabilities are non-negative and sum to 1). However, from Eq.~\ref{eq:op_policy}, it is evident that Eq.~\ref{eq:lag_p} represents the partition function of the optimal policy derived from the reward function $r(x, y)$. The essential insight behind the DPO algorithm is that it enforces specific constraints on the otherwise underdetermined Plackett-Luce family (particularly the Bradley-Terry model) of preference models. This ensures retention of the class of reward functions that can be represented while explicitly rendering the optimal policy in Eq.~\ref{eq:op_policy} analytically computable for every prompt $x$.

\subsection{Instability of Actor-Critic Algorithms} Our framework also facilitates diagnosing instabilities observed in conventional actor-critic methods utilized in RLHF, such as PPO. We follow the RLHF workflow, concentrating on the RL fine-tuning stage described in Section~\ref{section:prelims}. We establish connections to the control-as-inference framework \cite{levine2018reinforcement} for the constrained RL problem presented in Eq.~\ref{eq:RL}. Suppose we have a parameterized policy $\pi_{\theta}(y\mid x)$ and aim to minimize the divergence $\mathbb{D}_{\text{KL}}[\pi_{\theta}(y\mid x) \mid\mid \pi^*(y\mid x)]$, where $\pi^*$ denotes the optimal policy from Eq.~\ref{eq:optimum_model} induced by the reward function $r_{\phi}(y, x)$. Through algebraic manipulation, this yields the optimization objective:

\begin{equation}\label{eq:AC}
    \max_{\pi_{\theta}}\mathbb{E}_{\pi_{\theta}(y\mid x)}\bigg[\underbrace{r_{\phi}(x, y) -\beta\log\sum_{y}\pi_\text{ref}(y\mid x)\exp\left(\frac{1}{\beta}r_{\phi}(x, y)\right)}_{f(r_{\phi}, \pi_\text{ref}, \beta)} - \underbrace{\beta\log\frac{\pi_{\theta}(y\mid x)}{\pi_\text{ref}(y\mid x)}}_{\text{KL}}\bigg]
\end{equation}

This objective is identical to the one optimized in previous studies 
\citep{ziegler2020finetuning, stiennon2022learning, bai2022training, ouyang2022training} which use the DPO-equivalent reward for the reward class $r_{\phi}$. In this context, the normalization term in $f(r_{\phi}, \pi_\text{ref}, \beta)$ can be viewed as the soft value function of the reference policy $\pi_\text{ref}$. Although this term does not influence the optimal solution, its absence can cause the policy gradient of the objective to exhibit high variance, leading to unstable learning. This normalization term can be incorporated by employing a learned value function, though this approach may also present optimization challenges. Alternatively, earlier works have normalized rewards by utilizing a human completion baseline, effectively a single-sample Monte-Carlo approximation of the normalization term. By contrast, the DPO reparameterization produces a reward function that eliminates the need for any baselines.

\section{Experiments}
In this section, we empirically assess DPO’s capability to train policies directly from preference data. Initially, in a controlled text-generation environment, we investigate how effectively DPO balances maximizing reward and minimizing KL-divergence relative to the reference policy, in comparison to standard preference learning methods such as PPO. Subsequently, we test DPO’s performance on larger models and more challenging RLHF tasks, including summarization and dialogue. Our findings show that, with minimal hyperparameter tuning, DPO generally matches or exceeds the performance of strong baselines like RLHF with PPO, and also outperforms methods that select the best of $N$ sampled trajectories according to a learned reward function. Prior to discussing these outcomes, we outline the experimental setup; further details are provided in Appendix~\ref{app:exp_details}.

\textbf{Tasks.} Our study investigates three distinct open-ended text generation challenges. Across all experiments, the algorithms learn a policy based on a preference dataset $\mathcal{D}=\bigl\{x^{(i)}, y_w^{(i)}, y_l^{(i)}\bigr\}_{i=1}^N$. In \textbf{controlled sentiment generation}, the input $x$ represents a movie review prefix from the IMDb dataset \cite{maas-EtAl:2011:ACL-HLT2011}, and the goal is for the policy to generate an output $y$ exhibiting positive sentiment. To facilitate a controlled assessment, we \textit{generate} preference pairs from model outputs using a pre-trained sentiment classifier that satisfies $p(\text{positive}\mid x,y_w)>p(\text{positive}\mid x,y_l)$. For supervised fine-tuning (SFT), GPT-2-large is fine-tuned until convergence on training reviews from the IMDb dataset (additional details in App~\ref{app:sentiment_details}). In the \textbf{summarization} task, $x$ corresponds to a Reddit forum post, and the policy is tasked with generating a summary $y$ capturing the key points of the post. Consistent with previous work, we utilize the Reddit TL;DR summarization dataset \citep{volske-etal-2017-tl} alongside human preference annotations collected by \citeauthor{stiennon2022learning}. The SFT model used is fine-tuned on human-written summaries of forum posts\footnote{\url{https://huggingface.co/CarperAI/openai_summarize_tldr_sft}} via the TRLX \citep{leandro_von_werra_2023_7790115} framework for RLHF. The human preference dataset was compiled by \citeauthor{stiennon2022learning} using samples generated by a different but comparably fine-tuned SFT model. Lastly, in the \textbf{single-turn dialogue} task, $x$ is a user query, which can range from astrophysics questions to relationship advice requests. The policy must generate an engaging and informative response $y$. We employ the Anthropic Helpful and Harmless dialogue dataset \citep{bai2022training}, which includes 170k human-automated assistant interactions. Each dialogue concludes with two model-generated responses by a large (though undisclosed) language model, accompanied by a human preference label indicating the favored response. In this scenario, as no pre-trained SFT model exists, we fine-tune an off-the-shelf language model solely on the preferred completions to establish the SFT baseline.

\textbf{Evaluation.} Our experimental setup employs two distinct evaluation strategies. To assess how well each algorithm optimizes the constrained reward maximization objective within the controlled sentiment generation context, we measure each algorithm’s performance by examining the frontier defined by the reward achieved and the KL-divergence from the reference policy; this frontier is computable since we have access to the true reward function (a sentiment classifier). Conversely, in real-world scenarios where the true reward function is unknown, we evaluate algorithms based on their \textit{win rate} relative to a baseline policy, utilizing GPT-4 as a stand-in for human judgment regarding summary quality and response helpfulness in the summarization and single-turn dialogue tasks, respectively. For summarization, the baseline consists of reference summaries from the test set; for dialogue, the baseline is the preferred response present in the test dataset. Although prior work suggests that language models can serve as superior automated evaluators compared to conventional metrics \citep{Chen2023ExploringTU}, we validate our use of GPT-4 for evaluation through a human study described in Sec.~\ref{sec:human-judgments}. Our findings reveal that GPT-4’s assessments strongly correlate with human judgments, and that human agreement with GPT-4 is generally comparable to or exceeds the level of agreement observed between human annotators themselves.

\textbf{Methods.} Beyond DPO, we assess several established methods for training language models to align with human preferences. At a basic level, we examine zero-shot prompting using \textbf{GPT-J} \citep{gpt-j} for summarization, and 2-shot prompting with \textbf{Pythia-2.8B} \citep{biderman2023pythia} in the dialogue task. We also evaluate the \textbf{SFT} model and the \textbf{Preferred-FT} model, which is fine-tuned via supervised learning on the chosen completion $y_w$ derived either from the SFT model (in controlled sentiment and summarization) or from a generic LM (in single-turn dialogue). Another pseudo-supervised technique is \textbf{Unlikelihood}~\citep{welleck2019neural}, which aims to increase the probability of $y_w$ while simultaneously \textit{decreasing} the probability of $y_l$; this method incorporates an optional coefficient $\alpha\in[0,1]$ to weight the unlikelihood term. We further investigate \textbf{PPO} \citep{schulman2017proximal} trained on a reward function inferred from preference data, as well as \textbf{PPO-GT}, an oracle that learns using the ground truth reward function accessible in the controlled sentiment setup. In the sentiment experiments, two PPO-GT implementations are used: an off-the-shelf version \cite{leandro_von_werra_2023_7790115} and a modified variant that normalizes rewards and fine-tunes hyperparameters to enhance performance (these adjustments are also applied when training conventional PPO with learned rewards). Lastly, we consider the \textbf{Best of $N$} baseline, which samples $N$ completions from the SFT model (or Preferred-FT for dialogue) and selects the highest-scoring response based on a reward function derived from the preference dataset. Although this method yields strong performance by separating reward model quality from PPO optimization, it is computationally expensive for even moderate values of $N$ due to the requirement of generating $N$ completions for every query at test time.

\subsection{How effectively can DPO optimize the RLHF objective?}

The KL-constrained reward maximization objective commonly employed in RLHF methods balances the exploitation of reward with limiting the policy from straying too far from the reference policy. Hence, when comparing different algorithms, it is important to consider both the reward obtained and the KL divergence; obtaining slightly higher reward at the cost of significantly higher KL divergence is not necessarily preferable. Figure~\ref{fig:frontier-tldr-main} illustrates the reward-KL frontier for several algorithms within the sentiment domain. We perform multiple training runs per algorithm, varying the hyperparameter controlling policy conservativeness for each run (target KL values in $\{3,6,9,12\}$ for PPO, $\beta \in \{0.05,0.1,1,5\}$, $\alpha \in \{0.05,0.1,0.5,1\}$ for unlikelihood, and random seeds for preferred-FT). This hyperparameter sweep totals 22 runs. After every 100 training steps until convergence, we evaluate each policy using a set of test prompts, calculating the average reward under the true reward function alongside the average sequence-level KL\footnote{That is, the sum of the KL divergences at each timestep.} relative to the reference policy $\text{KL}\left(\pi \mid \mid \pi_\text{ref}\right)$. Our findings show that DPO delivers by far the most efficient frontier, attaining the highest reward while maintaining low KL divergence. This outcome is significant for several reasons. First, although DPO and PPO optimize the same objective, DPO is markedly more efficient; its reward/KL tradeoff strictly surpasses that of PPO. Second, DPO achieves a superior frontier compared to PPO, \emph{even in cases where PPO has access to ground truth rewards} (PPO-GT).

\subsection{Is DPO capable of scaling to real preference datasets?}
\label{sec:dpo-real-datasets}
We proceed to assess the fine-tuning effectiveness of DPO on tasks involving summarization and single-turn dialogue. In the context of summarization, automatic metrics like ROUGE often show weak correlation with human preferences~\citep{stiennon2022learning}, and earlier research has demonstrated that fine-tuning language models via PPO using human preference data yields better summaries. We compare various methods by generating completions on the TL;DR summarization dataset's test split and calculating the average win rate against reference completions in that test set. For all methods, completions are sampled across a temperature range from 0.0 to 1.0, with win rates depicted in Figure~\ref{fig:frontier-tldr-main} (right). DPO, PPO, and Preferred-FT all fine-tune the identical GPT-J SFT model\footnote{\url{https://huggingface.co/CarperAI/openai_summarize_tldr_sft}}. Our findings indicate that DPO achieves a win rate close to 61\% at temperature 0.0, outperforming PPO’s approximately 57\% win rate at its best sampling temperature of 0.0. Additionally, DPO attains a higher peak win rate compared to the best of $N$ baseline. It is important to mention that we did not extensively tune DPO’s $\beta$ hyperparameter, suggesting these results might underestimate DPO’s full capabilities. Furthermore, we observe that DPO is significantly more stable with respect to sampling temperature than PPO, whose performance can fall back to the base GPT-J model’s level at elevated temperatures. Preferred-FT does not show notable improvement over the SFT model. A direct comparison between DPO and PPO in human evaluations is presented in Section~\ref{sec:human-judgments}, where DPO samples drawn at temperature 0.25 were favored 58\% of the time over PPO samples at temperature 0.

For single-turn dialogue, we assess various approaches on a subset of the test portion of the Anthropic HH dataset \citep{bai2022training} that involves a single step of human-assistant interaction. GPT-4 evaluations utilize the preferred completions from the test set as a reference to calculate the win rate for each method. Since there is no established SFT model for this task, we begin with a pre-trained Pythia-2.8B model, apply Preferred-FT to train a reference model on the selected completions to ensure the completions remain within the model’s distribution, and subsequently train with DPO. We also benchmark against the top completion from 128 Preferred-FT samples (noting that the Best of $N$ baseline plateaus at 128 completions for this task; see Appendix Figure~\ref{fig:best-of-n}) as well as a 2-shot prompted version of the Pythia-2.8B base model. We observe that DPO performs comparably or better at the optimal temperatures for each method. Additionally, we evaluate an RLHF model trained via PPO on the Anthropic HH dataset \footnote{\url{https://huggingface.co/reciprocate/ppo_hh_pythia-6B}} from a recognized source \footnote{\url{https://github.com/CarperAI/trlx/tree/main/examples/hh}}, but fail to identify any prompt or sampling temperature that surpasses the base Pythia-2.8B model’s performance. Considering our TL;DR findings and that both methods optimize the same reward function, we treat Best of 128 as a rough estimate of PPO-level performance. Overall, DPO stands out as the only computationally efficient approach that improves upon the preferred completions in the Anthropic HH dataset, offering comparable or superior results relative to the computationally intensive Best of 128 baseline. Lastly, Figure~\ref{fig:dialogue-main} demonstrates that DPO attains its peak performance relatively rapidly.

\subsection{Extension to a novel input distribution}

To further assess the robustness of PPO and DPO when faced with distribution shifts, we test the PPO and DPO policies derived from our Reddit TL;DR summarization experiment on a different data distribution, specifically news articles from the test set of the CNN/DailyMail dataset \citep{nallapati-etal-2016-abstractive}, utilizing the optimal sampling temperatures identified from TL;DR (0 and 0.25). The outcomes are shown in Table~\ref{tab:ood}. We calculated the GPT-4 win rate by comparing against the ground-truth summaries in these datasets, employing the same GPT-4 (C) prompt used for Reddit TL;DR, but substituting the phrase ``forum post'' with ``news article''. Under this new distribution, DPO continues to significantly outperform PPO. This result offers preliminary evidence that DPO policies can generalize at least as well as PPO policies, despite DPO not leveraging the additional unlabeled Reddit TL;DR prompts that PPO utilizes.

\subsection{Corroborating GPT-4 evaluations with human judgments}
\label{sec:human-judgments}
We perform a human evaluation study to confirm the reliability of GPT-4’s assessments, using data from the TL;DR summarization experiment and two distinct GPT-4 prompts. The \textbf{GPT-4 (S)} (simple) prompt requests which summary better captures the important information in the post. The \textbf{GPT-4 (C)} (concise) prompt additionally asks which summary is more concise; this prompt is used because we observe that GPT-4 tends to favor longer, more repetitive summaries with the \textbf{GPT-4 (S)} prompt, unlike human preferences. The full prompts can be found in Appendix~\ref{app:prompts}. We conduct three pairwise comparisons, selecting the highest-performing method (DPO, temp. 0.25), the lowest-performing method (PPO, temp. 1.0), and a mid-level method (SFT, temp. 0.25) to encompass a range of summary quality; all are compared against the greedily-sampled PPO (its best temperature setting). Our findings indicate that for both prompts, GPT-4 agrees with humans roughly as often as humans agree amongst themselves, indicating that GPT-4 is a valid surrogate for human evaluation (due to the limited number of human raters, multiple human judgments were collected only for the DPO and PPO-1 comparisons). Overall, the \textbf{GPT-4 (C)} prompt tends to produce win rates more aligned with human judgments; accordingly, we use this prompt for the primary results in Section~\ref{sec:dpo-real-datasets}. Additional information about the human study, including the web interface shown to raters and the list of volunteers, is available in Appendix~\ref{app:human-study}.

\section{Discussion}
Learning from preferences offers a robust and scalable approach to developing capable and aligned language models. We have presented DPO, a straightforward training method that enables language models to be trained from preferences without relying on reinforcement learning. Instead of forcing the preference learning task into a conventional RL framework to utilize standard RL algorithms, DPO establishes a direct connection between language model policies and reward functions, allowing the model to be trained to meet human preferences \textit{directly} using a simple cross-entropy loss, without the need for reinforcement learning or compromising generality. With minimal hyperparameter tuning, DPO achieves performance on par with or better than existing RLHF techniques, including those based on PPO; thus, DPO significantly lowers the barriers for training language models from human preferences.

\textbf{Limitations \& Future Work.} Our findings open up several important avenues for further investigation. How well does the DPO policy generalize to out-of-distribution data compared to approaches relying on explicit reward functions? Preliminary evidence suggests that DPO policies generalize comparably to PPO-based models, but more extensive research is required. For instance, can self-labeling with the DPO policy similarly leverage unlabeled prompts effectively? Moreover, how does reward over-optimization appear within the direct preference optimization framework, and could the slight performance drop observed in Figure~\ref{fig:dialogue-main}-right be an example of this? Although we assess models up to 6B parameters, scaling DPO to cutting-edge models with much larger sizes represents an exciting path for future exploration. Concerning evaluations, we observe that GPT-4’s win rate measurements are influenced by the prompt used; subsequent work may focus on optimizing methods to extract high-quality assessments from automated systems. Lastly, DPO’s potential extends beyond training language models from human preferences and includes applications in training generative models across other modalities.

\end{document}